    \item Consider the following compression scheme for binary sequences.
We divide the binary sequences into blocks of size 16. For each
block if the number of zeros is greater than or equal to 8 we store a
0, otherwise we store a 1. If the sequence is random with probability of zero 0.9, compute the rate and average distortion (Hamming metric). 

% Problem 4: Shannon lower bound for the rate distortion function
%%%%%%%%%%%%%%%%%%%%%%%
\vskip .3in
\item{\bf Shannon lower bound for the rate distortion function.}\\
 Consider a source $X$ with a distortion measure $d(x,\xh)$ that satisfies the following property: all columns of the distortion matrix are permutations of the set $\{d_1,d_2,\ldots,d_m\}$.  Define the function
\eq
\phi(D) = \max_{\pu: \sum_{i=1}^m p_i d_i \leq D} H(\pu).
\en
The Shannon lower bound on the rate distortion function is proved by the following steps:
\begin{enumerate}
\item Show that $\phi(D)$ is a concave function of $D$.
\item Justify the following series of inequalities for $I(X;\Xh)$ if $\E d(X,\Xh) \leq D$, 
\ea
I(X;\Xh) &=& H(X) - H(X|\Xh) \\
&=& H(X) - \sum_{\xh} p(\xh) H(X|\Xh = \xh) \\
&\geq& H(X) - \sum_{\xh} p(\xh) \phi( D_{\xh}) \\
&\geq& H(X) - \phi\left( \sum_{\xh} p(\xh) D_{\xh} \right) \\
&\geq& H(X) -\phi(D),
\an
where $D_{\xh} = \sum_x p(x|\xh) d(x,\xh)$.   

\item Argue that
\eq
R(D) \geq H(X) -\phi(D),
\en
which is the Shannon lower bound on the rate distortion function.

\item If in addition, we assume that the source has a uniform distribution and that the rows of the distortion matrix are permutations of each other, then  $R(D) = H(X) - \phi(D)$, i.e., the lower bound is tight.
\end{enumerate}

% Solution of Problem 4: Shannon lower bound for the rate distortion function
%%%%%%%%%%%%%%%%%%%%%%%
\vskip .2in
\begin{solution}{\bf Solution: Shannon lower bound on the rate distortion function.}\\
\begin{enumerate}
\item We define
\eq
\phi(D) = \max_{\pu: \sum_{i=1}^m p_i d_i \leq D} H(\pu).
\en
From the definition, if $D_1 \geq D_2$, then $\phi(D_1) \geq \phi(D_2)$ since the maximization is over a larger set.  Hence $\phi(D)$ is a monotonic increasing function.

To prove concavity of $\phi(D)$, consider two levels of distortion $D_1$ and $D_2$ and let $\pu^{(1)}$ and $\pu^{(2)}$ achieve the maxima in the definition of $\phi(D_1)$ and $\phi(D_2)$.  Let $\pu^{(\l)}$ be the mixture of the two distributions, i.e.,
\eq
\pu^{(\l)} = \l \pu^{(1)} + (1-\l) \pu^{(2)}.
\en
Then the distortion is a mixture of the two distortions
\eq
D_{\l} = \sum_i p_i^{(\l)} d_i = \l D_1 + (1-\l) D_2.
\en
Since entropy is a concave function, we have
\eq
H(\pu^{(\l)}) \geq \l H(\pu^{(1)}) + (1-\l) H(\pu^{(2)}).
\en
Hence
\ea
\phi(D_{\l}) &=& \max_{\pu : \sum p_i d_i = D_{\l}} H(\pu) \\
&\geq& H(\pu^{(\l)}) \\
&\geq& \l H(\pu^{(1)}) +(1-\l) H(\pu^{(2)}) \\
&=& \l \phi(D_1) + (1-\l) \phi(D_2),
\an
proving that $\phi(D)$ is a concave function of $D$.

\item For any $(X,\Xh)$ that satisfy the distortion constraint, we have
\ea
I(X;\Xh) &\stackrel{(a)}{=}& H(X) - H(X|\Xh) \\
&\stackrel{(b)}{=}& H(X) - \sum_{\xh} p(\xh) H(X|\Xh = \xh) \\
&\stackrel{(c)}{\geq}& H(X) - \sum_{\xh} p(\xh) \phi( D_{\xh}) \\
&\stackrel{(d)}{\geq}& H(X) - \phi( \sum_{\xh} p(\xh) D_{\xh} ) \\
&\stackrel{(e)}{\geq}& H(X) -\phi(D),
\an
where \\
(a) follows from the definition of mutual information, \\
(b) from the definition of conditional entropy, \\
(c) follows from the definition of $\phi(D_{\xh})$ where $D_{\xh} = \sum p(x|\xh) d(x,\xh) = \sum p(x|\xh) d_{i'}$ that $H(p(x|\xh)) \leq \phi(D_{\xh}) $ \\
(d) follows from Jensen's inequality and the concavity of $\phi$, and \\
(e) follows from the monotonicity of $\phi$ and the fact that $\sum
p(\xh) D_{\xh} = \sum p(x,\xh) d(x,\xh) \leq D$. \\
Hence, from the definition of the rate distortion function, we have
\ea
R(D) &=& \min_{p(\xh|x) : \sum p(x,\xh) d(x,\xh) \leq D} I( X;\Xh) \\
&\geq& H(X) - \phi(D),
\an
which is the Shannon lower bound on the rate distortion function.

\item Let $\pu^* = (p^*_1,p^*_2,\ldots, p^*_m)$ be the distribution that achieves the maximum in the definition of the $\phi(D)$.  Assume that the source has a uniform distribution and that the rows of the distortion matrix are permutations of each other.  Let the distortion matrix be $[a_{ij}]$.  We can then choose $p(\xh)$ to have a uniform distribution and  choose $p(x=i|\xh=j) = p^*_k$, if $a_{ij} = d_k$. For this joint distribution,
\ea
p_x(i)  &=& \sum_j p_{\xh}(j) p_{x|\xh} (i|j) \\
&=& \sum_j {\frac{1}{m}} p^*_k \\
&=& {1\over m}
\an
since the rows of the distortion matrix are permutations of each other and therefore each element $p^*_k$, $k=1,2,\ldots, m$ occurs once in the above sum. Hence the distribution of $x$ has the desired source distribution.  For this joint distribution, we have
\ea
\sum_{i,j} p_{x,\xh}(i,j) a_{ij} 
& =& \sum_j {1\over m} \sum_i p_{x|\xh}(i|j) a_{ij} \\
& =& \sum_j {1\over m} \sum_k p^*_k d_k \\
&=& \sum_j {1\over m} D \\
&=& D,
\an
the desired distortion.   The mutual information 
\ea
I(X;\Xh) &=& H(X) - H(X|\Xh) \\
&=& H(X) - \sum_j {1\over m} H(X|\Xh= j) \\
&=& H(X) - \sum_j {1\over m} H(\pu^*) \\
&=& H(X) - \sum_j {1\over m} \phi(D) \\
&=& H(X) - \phi(D).
\an
Hence using this joint distribution in the definition of the rate distortion function
\ea
R(D) &=& \min_{p(\xh|x) : \sum p(x,\xh) d(x,\xh) \leq D} I( X;\Xh) \\
&\leq& I(X;\Xh) \\
&=& H(X) -\phi(D).
\an
Combining this with the Shannon lower bound on the rate distortion function, we must have equality in the above equation and hence we have equality in the Shannon lower bound.
\end{enumerate}
\end{solution}